% ----------------------
\subsection{Section 47 (Circa 8 March 1831)}
% ----------------------
	\label{DC47}
	
	% -----
	\subsubsection{Historical Background}
	% -----
		\label{DC47-HB}
		
		% --
		\paragraph{JSP}
		% --
		
			``This revelation appointed John Whitmer to take up the work formerly done by Oliver Cowdery as scribe and recorder---and added a new responsibility. Cowdery, who was `called...to write for' JS by April 1829, was the principal scribe for the Book of Mormon and JS's early revelations and, once the Church of Christ was organized in April 1830, kept minutes at meetings. He also served as scribe for early portions of JS's inspired revision of the Bible. By the time of this revelation, however, Cowdery was away on a mission, and Whitmer had assumed some of Cowdery's formal duties. He succeeded Cowdery as scribe for JS's Bible revision and also assisted JS in gathering and copying revelations, a work that culminated in the creation of a manuscript book of revelations. The following revelation not only formalized Whitmer's assignment to `assist my servent Joseph' in such writing duties but added that he should also `write \& keep a regulal [regular] history.' The revelation then reiterated this dual appointment, directing that Whitmer `Keep the Church Record \& History continually for Oliver I have appointed to an other office.'
			
			``Before the establishment of the church, JS's primary need for scribal assistance was for recording revelatory texts, as demonstrated by Cowdery's work on the translation of the Book of Mormon as well as his careful recording of at least sixteen revelations. The need for other kinds of record keeping increased once the Church of Christ was organized. The first lines of a revelation recorded on 6 April 1830, the day of organization, proclaimed that `there Shall a Record be kept among you,' and the church's foundational `Articles and Covenants' commanded that a record of members be kept and that members be provided with certificates of good standing that could be presented to the various branches of the church. Cowdery responded by keeping these records until he departed in October 1830 for a mission to the American Indians. David Whitmer was temporarily appointed to record conference minutes, and in March 1831, John Whitmer was appointed to replace Cowdery. About the time this revelation was dictated, John Whitmer began inscribing Revelation Book 1.
			
			``John Whitmer was a willing scribe but a reluctant historian. At the 9 April 1831 conference he responded affirmatively by keeping the minutes when `appointed to keep the Church record \& history by the voice of ten Elders.' But according to Whitmer's own history, when JS told him that in addition to keeping church records `you must also keep the Church history,' he initially declined: `I would rather not do it but observed that the will of the Lord be done, and if he desires it, I desire that he would manifest it through Joseph the Seer.' As a result of this conversation, JS dictated the revelation presented here. Whitmer likely composed its heading when copying it into Revelation Book 1, acknowledging that the text was dictated `in consequenc[e] of not feeling reconsiled to write at the request of Joseph with[o]ut a commandment.'
			
			``Prior to the dictation of this revelation, record keeping in the church was often sporadic and incomplete. After Whitmer was appointed church historian and recorder, the number of documents recording church history increased substantially. Minutes of church conferences generally contained more detail than they had previously, and Whitmer's creation of Revelation Book 1 preserved most of JS's early revelations. Whitmer also wrote a ninety-six-page narrative history that primarily described events from fall 1830 through the mid-1830s.''
		
	% -----
	\subsubsection{Versions}
	% -----
		\label{DC47-V}
		
		See the \href{https://drive.google.com/file/d/1goAvNz8jS4R8slYfMG_db55Ty2z_vmgK/view?usp=sharing}{sections comparison} document.
	
		\begin{compactdesc}
			\item[JSP] \hyperref[EMS-1831-JS-CIR-8-MAR-1831-B]{Circa 8 March 1831}
			\item[BC] Chapter 50, 114
			\item[DC 1835] Section 63, 190
			\item[TS] 5:2 (15 January 1844), 401
			\item[MS] 5:6 (November 1844), 82
			\item[DC 1844] Section 63, 288
		\end{compactdesc}

	% -----
	\subsubsection{Section 47:1} % *DC47-1 
	% -----
		\label{DC47-1}

		BEHOLD, it is expedient in me that my servant John should write and keep a regular history, and assist you, my servant Joseph, in transcribing all things which shall be given you, until he is called to further duties.

		% \begin{framed}
		% \scriptsize
			% \begin{compactdesc}
				% \item[JSP] 
				% % \item[BC] 
				% % \item[]
			% \end{compactdesc}
		% \end{framed}

	% -----
	\subsubsection{Section 47:2} % *DC47-2 
	% -----
		\label{DC47-2}

		Again, verily I say unto you that he can also lift up his voice in meetings, whenever it shall be expedient.

		% \begin{framed}
		% \scriptsize
			% \begin{compactdesc}
				% \item[JSP] 
				% % \item[BC] 
				% % \item[]
			% \end{compactdesc}
		% \end{framed}

	% -----
	\subsubsection{Section 47:3} % *DC47-3 
	% -----
		\label{DC47-3}

		And again, I say unto you that it shall be appointed unto him to keep the church record and history continually; for Oliver Cowdery I have appointed to another office.

		% \begin{framed}
		% \scriptsize
			% \begin{compactdesc}
				% \item[JSP] 
				% % \item[BC] 
				% % \item[]
			% \end{compactdesc}
		% \end{framed}

	% -----
	\subsubsection{Section 47:4} % *DC47-4 
	% -----
		\label{DC47-4}

		Wherefore, it shall be given him, inasmuch as he is faithful, by the \hyperref[COMFORTER]{Comforter}, to write these things.  Even so.  Amen.

		% \begin{framed}
		% \scriptsize
			% \begin{compactdesc}
				% \item[JSP] 
				% % \item[BC] 
				% % \item[]
			% \end{compactdesc}
		% \end{framed}